% modeling/part_intro.tex

\chapter*{Modeling}
\phantomsection
\addcontentsline{toc}{chapter}{Modeling}

% ============================================================
\section*{What do we already know?}
% ============================================================

\begin{remark}
The previous parts established the conceptual foundations of AIM.
\end{remark}

\begin{remark}
Intrinsic geometry, railway states, runtime evaluation, realization,
and reality-aware interpretation have been developed as coherent
theoretical constructs.
\end{remark}

\begin{remark}
These concepts define what railway alignments are, how they evolve,
and how they interact with the physical world.
\end{remark}

\begin{remark}
However, concepts alone do not yet define how alignment information is
represented within an operational information model.
\end{remark}

% ============================================================
\section*{What is still missing?}
% ============================================================

\begin{remark}
An information model requires explicit internal structures.
\end{remark}

\begin{remark}
Geometric concepts must be transformed into representations that can be
stored, exchanged, reconstructed, analyzed, and optimized.
\end{remark}

\begin{remark}
Traditional railway systems frequently rely on coordinate lists,
geometric primitives, or implementation-specific data structures.
\end{remark}

\begin{remark}
While such representations may be sufficient for storage purposes, they
often obscure the intrinsic structure of the alignment and limit
subsequent analysis.
\end{remark}

\begin{remark}
The missing element is therefore a modeling framework capable of
capturing railway semantics while remaining compact, interpretable,
and computationally efficient.
\end{remark}

% ============================================================
\section*{What comes next?}
% ============================================================

\begin{remark}
This part introduces the modeling layer of AIM.
\end{remark}

\begin{remark}
The central objective is the construction of sparse, semantic alignment
models.
\end{remark}

\begin{remark}
Within AIM, an alignment is not primarily represented as sampled
coordinates.

Instead, it is represented through a compact collection of geometric
states, transition descriptions, and semantic relationships from which
geometry can be reconstructed whenever required.
\end{remark}

\begin{remark}
The chapters that follow introduce:
\begin{itemize}
\item sparse alignment representations,
\item transition modeling,
\item transition classification,
\item higher-level alignment semantics,
\item and Schwerpunkt-based modeling concepts.
\end{itemize}
\end{remark}

\begin{remark}
These structures form the operational core of alignment-based
information modelling.
\end{remark}

% ============================================================
\section*{Relation to the AIM Roadmap}
% ============================================================

\begin{remark}
Within the AIM roadmap, the present part corresponds to the transition
from conceptual railway descriptions toward explicit alignment models.
\end{remark}

\begin{center}
\begin{tikzpicture}[
  node distance=1.4cm,
  box/.style={
    draw,
    rounded corners,
    minimum width=4.8cm,
    minimum height=0.9cm,
    align=center
  }
]

\node[box] (concepts)
{Concepts};

\node[box, below=of concepts] (sparse)
{Sparse Model};

\node[box, below=of sparse] (alignment)
{Alignment Model};

\draw[->, thick] (concepts) -- (sparse);
\draw[->, thick] (sparse) -- (alignment);

\end{tikzpicture}
\end{center}

\begin{remark}
The progression
\[
\text{Concepts}
\rightarrow
\text{Sparse Model}
\rightarrow
\text{Alignment Model}
\]
marks the transition from theoretical railway descriptions toward
operational information structures.
\end{remark}

\begin{remark}
These models subsequently become the basis for runtime services,
analysis pipelines, optimization procedures, and engineering
applications.
\end{remark}

% ============================================================
\begin{principle}[Sparse Modeling Principle]
\label{princ:sparse-modeling}
Within AIM, railway alignments should be represented through compact
semantic structures from which geometric, operational, and realization-
aware information can be reconstructed when required.
\end{principle}

% ============================================================
\begin{summary}
Concepts explain railway structure.

Reality explains interaction with the physical world.

Modeling explains how railway knowledge is represented.

The present part therefore establishes the sparse and semantic
information structures that form the core of alignment-based
information modelling.
\end{summary}
